\documentclass[10pt,a4paper]{article}
\usepackage[T1]{fontenc}
\usepackage[margin=22mm]{geometry}
\usepackage{lmodern,amsmath,amssymb,amsthm,mathtools,booktabs,microtype}
\usepackage[hidelinks]{hyperref}
\newtheorem{theorem}{Theorem}
\title{A hard-prime ES shell missed by first-degree moments:\newline exact counterexample, positive repairs and a progression}
\author{The Clankers}
\date{14 September 2026 --- research preprint}
\begin{document}\maketitle
\begin{abstract}
An occupied original divisor shell at a prime congruent to one modulo 840 has
negative-definite aggregate first-degree signed localizing form. We give its
complete finite residue calculation, deterministic primality certificate, and
unique marked witness. Retaining the exterior channel separately or applying
a discrete quartic both certify positivity. The witness extends to an explicit
all-integer arithmetic progression with infinitely many prime members. This
is not an ES counterexample, new prime coverage, or a claim of historical priority.
\end{abstract}

\paragraph{Original definitions.}
For $p/4<a<p/2$ and $R=4a-p$, each divisor $u\mid a^2$ has two channel labels,
with $g_E=4u+1$, $g_M=u+a$ and $D=R/\gcd(R,g_c)$. The original rational
triples are
\[
 E:\left(a,\frac{pa+a^2/u}R,\frac{pa+p^2u}R\right),\qquad
 M:\left(a,\frac{p(a+a^2/u)}R,\frac{p(a+u)}R\right).
\]
They sum reciprocally to $4/p$ since $(RY-pa)(RZ-pa)=p^2a^2$.
The two nonfirst reduced denominators are both $D$, using $4a\equiv p\pmod R$
and $\gcd(pa,R)=1$. Thus $D=1$ is exactly a hit, and all failed $D$ are odd
and at least three. These are the inherited Type I/II coordinates \cite{ET}.

Let $\mu=\sum_{c,u}\delta_{1/D}$ retain every original state. Define
\[
 L_1(\mu)=\max_{\deg P\le1,\ P(1)=1}
 \frac32\int(x-1/3)P(x)^2\,d\mu(x).
 \tag{1}
\]
Successful atoms contribute one and failed atoms nonpositively, so $L_1\le H$,
where $H$ is the original hit count. The preceding finite screen found no
occupied shell with $L_1\le0$ through prime parameter 50000; its source
explicitly made no universal converse claim \cite{Prior}.

\begin{theorem}[An occupied hard shell with no aggregate first-degree certificate]
Let
\[
 p=19748674681\equiv1\pmod{840},\qquad
 a=4937168704=2^6\cdot13^4\cdot37\cdot73,\qquad R=135.
 \tag{2}
\]
Then $p$ is prime, the shell contains exactly one original hit, but
\[
 L_1(E+M)=-\frac{50887012}{144801015}<0.
 \tag{3}
\]
\end{theorem}
\begin{proof}
A full order certificate uses base 11 and
$p-1=2^3\cdot3\cdot5\cdot7\cdot23510327$.
Every factor is prime by trial division, the largest requiring divisors only
through 4848. Direct modular powering gives $11^{p-1}\equiv1$ and the residues
\[
\begin{array}{c|rrrrr}
q&2&3&5&7&23510327\\\hline
11^{(p-1)/q}\bmod p&19748674680&10039250736&3428156451&13907516733&12674160518.
\end{array}
\]
The gcd with $p$ of every displayed residue minus one is one. Thus modulo any
prime divisor of $p$ the order of 11 is divisible by $p-1$, forcing that divisor
to be at least $p$. This proves primality. The verifier also independently
tries every odd divisor of $p$ through 140529.

In the finite group algebra of $(\mathbb Z/135\mathbb Z)^\times$, form
\[
 (\sum_{i=0}^{12}[2]^i)(\sum_{j=0}^{8}[13]^j)
 (\sum_{k=0}^{2}[37]^k)(\sum_{l=0}^{2}[73]^l).
 \tag{4}
\]
Applying the two gates to every residue with its coefficient gives the
following complete table. The factors sum to 1053 actual divisors; direct
integer divisor enumeration gives the same table independently.

\newpage
\noindent\emph{Proof continued.}
\[
\begin{array}{c|rrrrrrrr}
D&1&3&5&9&15&27&45&135\\\hline
E&1&39&53&80&69&144&244&423\\
M&0&40&54&80&68&144&244&423\\
E+M&1&79&107&160&137&288&488&846
\end{array}\tag{5}
\]
This finite multiplication is itself an executable certificate: start with
coefficient one at the group identity; for each factor in (4), replace the
coefficient array by the sum of its displayed translates. No generated-group
support or unlimited multiplicity substitutes for the exponents.

Put $S_j=\sum_{c,u}\gcd(R,g_c)^j$ and $A_i=3S_{i+1}-RS_i$. The first five
moments from (5) are
\[
(S_0,S_1,S_2,S_3,S_4)=(2106,13962,315738,12455226,722183418).
\]
They yield
\[
(A_0,A_1,A_2)=(-242424,-937656,-5258952),\qquad
A_0A_2-A_1^2=395697405312>0.
\]
The signed form on linear polynomial coefficients is a positive scalar times
a diagonal congruence of $\begin{psmallmatrix}A_0&A_1\\A_1&A_2\end{psmallmatrix}$.
It is negative definite. Substituting $P=1+b(x-1)$ and maximizing the resulting
quadratic in $b$ gives (3). Table (5) proves $H=1$.
\end{proof}

\paragraph{Positive repair one: keep the channel.}
For a genuine disjoint partition of original states $\mathcal A=\bigsqcup A_i$,
\[
 \boxed{\sum_i\max\{0,\lceil L_d(\mu_i)\rceil\}\le H.}\tag{6}
\]
It strengthens under refinement. Indeed each summand bounds its own integer
hit count, and every coarse normalized polynomial is available on each finer
piece. Therefore $L_{\rm coarse}\le\sum L_{\rm fine}$; the ceiling inequality
and discarding negative summands prove monotonicity.

On the exterior channel of (2), the fixed polynomial $P=(36x-1)/35$ gives
\[
 \boxed{\frac32\sum_E(x-1/3)P(x)^2=\frac{64376}{196875}>0.}\tag{7}
\]
The full coefficient form is $\operatorname{diag}(J_E,J_M)$. The aggregate
form is its restriction along $P\mapsto(P,P)$, with pullback $J_E+J_M$.
The restriction is not an isomorphism; (7) remains a direction of the original
full source even though that diagonal restriction is negative definite.
These are disjoint original labels. After other observation maps, the actual
mixed Gram must still be transported rather than presumed block diagonal.

\paragraph{Positive repair two: use the discrete spectrum.}
Define
\[
 \Phi_j(x)=x\prod_{\nu=1}^j\frac{(2\nu+1)x-1}{2\nu}.
\]
At $D=2k+1>1$, this is zero if $k\le j$ and otherwise equals
$(-1)^j\binom{k-1}{j}/D^{j+1}$. Also $\Phi_j(1)=1$.
Thus odd $j$ bounds the original hit indicator below and even $j$ bounds it
above. From (5), exact rational evaluation gives
\[
 \boxed{\sum_{E+M}\Phi_3(1/D)=\frac{1297691}{2460375}>0,\qquad
 \sum_{E+M}\Phi_4(1/D)=\frac{115374406}{110716875}.}\tag{8}
\]
The bracket contains exactly the integer one. Since
$48\Phi_3=105x^4-71x^3+15x^2-x$, positivity follows from the explicit integer
inequality $105S_4-71RS_3+15R^2S_2-R^3S_1>0$.
The companion note proves factorial uniform errors for the entire filter family.

\newpage
\paragraph{The actual witness and its inverse.}
The sole hit is
\[
 (c,u,h,r,s,\kappa)=(E,141879296,26,2336,81289,341725215823),
\]
with original centered exponent vector $(5,-3,-1,1)$ and raw denominators
\[
 \boxed{(4937168704,\ 722241027794932022,\
 409884191461517484986392768).}\tag{9}
\]
Directly $a=hrs$, $u=hr^2$, and $135\kappa=pr+s$.
Clearing the reciprocal sum of $(hrs,hs\kappa,phr\kappa)$ gives
$(p+135)\kappa=4hrs\kappa$, proving ES. Its inverse is
$u=a^2/(135Y-pa)$, followed by
$(h,r,s)=(d_0^2/u,u/d_0,a/d_0)$ with $d_0=\gcd(a,u)$.
The slope $2336/81289<1/8$ lies in the programme's lower positive chamber.
No denominator permutation or exponent occurrence has been discarded.

\begin{theorem}[All-integer progression with a canonical marked return]
For every integer $j\ge0$, set
\[
 t=73+945j,\qquad C=67632448,\qquad
 n=4Ct-135=19748674681+255650653440j,
\]
\[
 h_0=26,\qquad r_0=32t,\qquad s_0=81289,\qquad a=Ct,
\]
\[
 \kappa_0=341725215823+8847406287328j+57265746370560j^2.
\]
Then the positive integers $(a,h_0s_0\kappa_0,nh_0r_0\kappa_0)$ solve ES at $n$.
Every $n$ is $1\pmod{840}$, and infinitely many members are prime.
\end{theorem}
\begin{proof}
Polynomial expansion proves $a=h_0r_0s_0$ and
$135\kappa_0=nr_0+s_0$. The same reciprocal calculation as for (9) proves the
identity; all terms are positive for $j\ge0$. The initial term and step are
coprime, so Dirichlet's theorem supplies infinitely many primes
\cite[Theorem 18.1]{Dirichlet}. No hypothesis that $t$ is prime is imposed.
If $d=\gcd(r_0,s_0)$, then $\gcd(d,135)=1$ because $\gcd(s_0,135)=1$.
The quotient equation gives $d\mid\kappa_0$. Therefore
\[
 (h,r,s,\kappa)=(h_0d^2,r_0/d,s_0/d,\kappa_0/d),\qquad u=h_0r_0^2
\]
is the canonical coprime marking and keeps all three ordered denominators.
On this image the parameter inverse is
$j=(n-19748674681)/255650653440$.
\end{proof}
The progression is an explicit specialization of the inherited selector.
Its existence is not claimed as a new general progression method, and its
later slopes are not all asserted to lie in the separate chamber.

\paragraph{Limits and reproducibility.}
The counterexample is to a possible universal \emph{converse} for the sufficient
first-degree moment test, not to any previously proved sufficient theorem.
It is not claimed least. Channel resolution itself is not asserted universally
sufficient: the full workbench gives another certified hard-prime shell at
which both degree-one channel maxima are negative and the same quartic is
positive. No fixed-degree universal sufficiency or ES proof is claimed.
The standalone standard-library verifier performs both complete histogram
calculations, checks the prime certificates and every labelled inverse, and
returns exact rational scores. No probabilistic primality assumption enters.

\begin{thebibliography}{3}\small
\bibitem{ET} Elsholtz, C., \& Tao, T. (2013). Counting the number of solutions
to the Erd\H{o}s--Straus equation on unit fractions. \emph{J. Aust. Math. Soc.,
94}(1), 50--105, \S2.
\bibitem{Prior} The Clankers (2026). \emph{Prime-power boundary kernels and
canonical ES count certificates}. ES workbench draft PR6, commit
\texttt{f234ed0c34100c39e52c6a4645793823e830a983}, equation M9 and finite screen.
\bibitem{Dirichlet} Sutherland, A. V. (2017). \emph{18.785 Lecture 18:
Dirichlet L-functions, primes in arithmetic progressions}, Theorem 18.1.
\end{thebibliography}
\end{document}
